\chapter{Modern Quantum Computing in 2026}
\label{ch:modern}

This chapter surveys the landscape a serious introductory student should know beyond foundational algorithms: hardware platforms, surface codes and logical qubits, the NISQ era, software ecosystems, fault-tolerance roadmaps, and post-quantum cryptography. Statements are tagged by evidence type; company roadmaps are not scientific facts and must be verified against peer-reviewed literature.

\section{Physical Qubit Platforms}
\label{sec:hardware}

\subsection{Leading modalities}

As of the \cutoff{} cutoff, several qubit modalities compete:
\begin{itemize}[noitemsep]
\item \textbf{Superconducting transmons:} microwave control, fast gates, planar fabrication; short $T_1,T_2$ relative to ions.
\item \textbf{Trapped ions:} long coherence, high fidelity, all-to-all connectivity in chains; slower gates.
\item \textbf{Neutral atoms:} reconfigurable arrays via optical tweezers; promising for error-correction layouts.
\item \textbf{Photonic:} room-temperature operation; challenges in deterministic two-qubit gates.
\item \textbf{Spin qubits:} semiconductor integration potential; scaling research active.
\end{itemize}

\subsection{Metrics that matter}

Compare platforms using gate fidelity ($>99.9\%$ two-qubit targets), coherence times ($T_1,T_2$), connectivity graph, and reproducible cross-calibrated benchmarks---not qubit count alone.

\begin{example}
A device advertising ``1000+ qubits'' with $10^{-2}$ two-qubit error rates cannot yet run fault-tolerant Shor at cryptographically relevant scales without extensive QEC overhead \cite{preskill2018}.
\end{example}

\section{Logical Qubits and the Surface Code}
\label{sec:surface-code}

\subsection{From physical to logical}

Fault-tolerant computation stores information in \emph{logical} qubits encoded across many \emph{physical} qubits. Syndrome extraction detects errors faster than they accumulate. The \emph{threshold theorem} guarantees arbitrary low logical error rates if physical errors per gate stay below a code-specific threshold.

\subsection{Surface code highlights}

The surface code \cite{acharya2024} arranges stabilizer checks on a 2D lattice with local interactions suitable for superconducting and neutral-atom layouts. Recent experiments demonstrate \emph{logical} error rates below \emph{physical} rates at increasing code distance---a milestone toward scalable QEC, though resource estimates for chemistry or RSA-scale factoring remain substantial.

\begin{example}
Distance-$d$ surface code requires $\approx d^2$ physical qubits per logical qubit and $d$ rounds of syndrome extraction per logical gate in rough estimates; exact constants depend on implementation and magic-state factories.
\end{example}

\section{Noise, Channels, and Mitigation}
\label{sec:nisq}

\subsection{NISQ definition}

NISQ devices \cite{preskill2018} have $50$--$500$ noisy qubits without full error correction. Useful tasks include variational algorithms, simulation prototypes, and benchmarking---not arbitrary-depth fault-tolerant programs.

\subsection{Mitigation versus correction}

Mitigation (Chapter~\ref{ch:qec}) reduces statistical bias in short circuits; correction suppresses logical error rates exponentially given sufficient redundancy. Long-term utility requires the latter for deep algorithms.

\section{Quantum Advantage and Benchmarks}
\label{sec:advantage}

\subsection{Random circuit sampling}

Random circuit sampling milestones \cite{arute2019} demonstrate quantum behavior difficult to simulate classically for \emph{specific} tasks. This is not universal speedup for everyday workloads.

\subsection{How to read claims}

Always ask: \emph{Which problem? Which metric? Which classical comparison?} Cross-entropy benchmarks, quantum volume, and application-specific demos measure different things.

\begin{example}
Sycamore-style supremacy argues classical simulation of a particular random circuit distribution is infeasible; it does not imply Gmail runs faster on quantum processors.
\end{example}

\section{Quantum Networking and Cryptography}
\label{sec:networking-crypto}

\subsection{Quantum key distribution}

QKD (Chapter~\ref{ch:protocols}) offers information-theoretic security under explicit device and channel trust models. Deployment exists in limited links; scaling requires quantum repeaters still under research.

\subsection{Post-quantum classical cryptography}

Shor's algorithm threatens RSA and elliptic-curve cryptography by efficiently solving hidden subgroup problems underlying those schemes. \emph{Post-quantum cryptography} (PQC) standardizes classical algorithms believed secure against quantum adversaries (lattice-based, hash-based, code-based). PQC deployment proceeds independently of quantum hardware timelines.

\section{Software Ecosystem (2026)}
\label{sec:software-ecosystem}

\subsection{Stack layers}

\begin{itemize}[noitemsep]
\item \textbf{Languages \& IR:} OpenQASM~3, Quil, vendor DSLs.
\item \textbf{SDKs:} Qiskit, Cirq, PennyLane, Braket, Azure Quantum, IBM Quantum Platform \cite{ibm2026}.
\item \textbf{Compilers:} routing, pulse-level optimization, error-aware transpilation.
\item \textbf{Simulators:} state-vector, density matrix, stabilizer, tensor-network.
\end{itemize}

\subsection{Hybrid workflows}

Near-term applications combine classical optimizers with quantum coprocessors (variational quantum eigensolvers, QAOA). HPC integration treats quantum devices as specialized accelerators with queueing and calibration schedules.

\section{Fault Tolerance Roadmaps}
\label{sec:ft-roadmaps}

\subsection{Magic states and distillation}

Non-Clifford gates (especially $T$) dominate resource estimates. \emph{Magic-state distillation} factories supply clean $T$ states to fault-tolerant circuits; their footprint often exceeds the computation proper.

\subsection{Company versus science}

Vendor announcements (e.g.\ processor codenames, qubit counts) are engineering milestones, not proofs of utility. Distinguish peer-reviewed logical-qubit demonstrations \cite{acharya2024} from marketing materials \cite{google2024}.

\section{Resource Estimation and Classical Simulation}
\label{sec:resources}

Exact simulation of $n$ qubits needs $O(2^n)$ complex amplitudes---infeasible beyond $\approx 50$ qubits on classical supercomputers for generic states. Clifford circuits simulate in polynomial time, clarifying why non-Clifford content matters for advantage claims.

\section{Open Questions at the \cutoff{} Cutoff}
\label{sec:open-questions}

\begin{itemize}
\item Fault-tolerant logical qubits at scale with practical space-time overheads.
\item Durable quantum advantage in chemistry, optimization, or machine learning.
\item Integration of quantum accelerators with classical HPC and secure PQC migration.
\item Standardized benchmarks correlating device metrics with application performance.
\end{itemize}

\section*{Further Reading}

See \cite{nielsen2010,preskill2018,ibm2026} and peer-reviewed surveys. Re-verify hardware claims against primary literature before citing in research.

\section*{Exercises}
\begin{enumerate}[label=\arabic*.]
\item Compare two hardware platforms on fidelity, coherence, and connectivity in a table.
\item Explain why surface codes favor 2D-local architectures.
\item Distinguish quantum advantage, quantum supremacy, and fault tolerance in one paragraph each.
\item Why does PQC migration not require a working quantum computer?
\item Estimate qubits needed to factor RSA-2048 using rough literature exponents (order-of-magnitude).
\item Name three software layers in the 2026 quantum stack.
\item What role do magic states play in fault-tolerant $T$ gates?
\item Critique a hypothetical headline: ``1000-qubit computer breaks encryption.''
\end{enumerate}
